%\documentclass[12pt,letterpaper]{article}



%%%
%% Required packages:
%%
\usepackage{etex}
\usepackage{amsmath}
\usepackage{amssymb}
\usepackage{amstext}
\usepackage{amsthm}
\usepackage{fancyhdr,fancybox}
\usepackage{mathrsfs} % need for \mathscr command
\usepackage{multicol}
\usepackage[normalem]{ulem}
%\usepackage{pictex,pictexplus}
% \usepackage{pictexwd}
\usepackage{rotating}
\usepackage{url}
\usepackage{xfrac}
\usepackage{booktabs}
\usepackage{color,soul}
\usepackage{comment}

\usepackage{hyperref}
\hypersetup{
    colorlinks=true,     % false: boxed links; true: colored links
    linkcolor=blue,      % color of internal links
    citecolor=blue,      % color of links to bibliography
    filecolor=blue,      % color of file links
    urlcolor=blue        % color of external links
}


\usepackage{tikz}
\usetikzlibrary{backgrounds,calc}

%%
%% Common formatting commands
%%
\setlength{\parindent}{0in}
\setlength{\textwidth}{7in}
\setlength{\evensidemargin}{-0.25in}
\setlength{\oddsidemargin}{-0.25in}
\setlength{\parskip}{.5\baselineskip}
\setlength{\topmargin}{-0.5in}
\setlength{\textheight}{9in}

%               Problem and Part
\newcounter{problemnumber}
\newcounter{partnumber}
\newcounter{subpartnumber}
\newcommand{\Problem}{%
  \refstepcounter{problemnumber}%
  \setcounter{partnumber}{0}%
  \setcounter{subpartnumber}{0}%
  \item[\makebox{\hfil\textbf{\theproblemnumber.}\hfil}]%
  }
\newcommand{\InProblem}[2][1.6]{%
  \refstepcounter{problemnumber}%
  \setcounter{partnumber}{0}%
  \setcounter{subpartnumber}{0}%
  \textbf{\theproblemnumber.}\ \ \parbox[t]{#1in}{#2}%
}
\newcommand{\InSmallProblem}[1]{\InProblem[1.05]{#1}}
\newcommand{\NoProblem}[1]{%
  \mbox{\hfil\textbf{#1}\hfil}%
  }
\newcommand{\UnProblem}{%
  \refstepcounter{problemnumber}%
  \setcounter{partnumber}{0}%
  \setcounter{subpartnumber}{0}%
  \mbox{\hfil\textbf{\theproblemnumber.}\hfil}%
  }
\newcommand{\InFlexProblem}[1]{%
  \refstepcounter{problemnumber}%
  \setcounter{partnumber}{0}%
  \setcounter{subpartnumber}{0}%
  \textbf{\theproblemnumber.}\ \ #1%
}
%%  Part:
\newcommand{\Part}{%
  \renewcommand{\thepartnumber}{(\alph{partnumber})}%
  \refstepcounter{partnumber}
  \setcounter{subpartnumber}{0}%
  \item[(\alph{partnumber})]%
}
\newcommand{\InPart}[2][1.75]{%
  \renewcommand{\thepartnumber}{(\alph{partnumber})}%
  \refstepcounter{partnumber}%
  \setcounter{subpartnumber}{0}%
  (\alph{partnumber})\ \ \parbox[t]{#1in}{#2}%
}
\newcommand{\UnPart}{%
  \refstepcounter{partnumber}%
  \renewcommand{\thepartnumber}{(\alph{partnumber})}%
  \setcounter{subpartnumber}{0}%
  (\alph{partnumber})%
}
\newcommand{\InLargePart}[1]{\InPart[3.05]{#1}}
\newcommand{\InFullPart}[1]{\InPart[6.2]{#1}}
\newcommand{\InSmallPart}[1]{\InPart[1.05]{#1}}
%% SubPart:
\newcommand{\SubPart}{%
  \refstepcounter{subpartnumber}%
  \item[(\roman{subpartnumber})]%
  }
\newcommand{\InSubPart}[2][2.25]{%
  \stepcounter{subpartnumber}%
  (\roman{subpartnumber})\ \ \parbox[t]{#1in}{#2}%
}
\newcommand{\InSmallSubPart}[1]{\InSubPart[1.5]{#1}}
\newcommand{\InTinySubPart}[1]{%
  \stepcounter{subpartnumber}%
  (\roman{subpartnumber})\ \ \makebox{#1}%
}
\newcommand{\InMySubPart}[2][2.25]{%
  \stepcounter{subpartnumber}%
  \parbox[t]{#1in}{\makebox[0.3in][l]{(\roman{subpartnumber})}\ \ {#2}}%
}
\newcommand{\TrueFalse}{\textbf{T}\ \ \textbf{F}\ \ }
\newcommand{\TFPart}[1]{% 
  \stepcounter{partnumber}%
  \setcounter{subpartnumber}{0}%
  \item[(\alph{partnumber})] \TrueFalse\parbox[t]{5.5in}{#1}}

\newcommand{\Points}[1]{(#1 points) \ }
\newcommand{\Point}[1]{(#1 point) \ }


%% For Answers:
\newcounter{answernumber}
\newcommand{\Answer}{%
  \refstepcounter{answernumber}%
  \setcounter{partnumber}{0}%
  \setcounter{subpartnumber}{0}%
  \item[\makebox{\hfil\textbf{\theanswernumber.}\hfil}]%
  }


% Setting up the headers for the first page
%\chead{} % will default to what is typed in the file.
\fancyhead[L]{\textsc{Math 22a}}
\fancyhead[R]{\textsc{Fall 2024}}
\fancyfoot[R]{
	\vspace*{\fill}\footnotesize\textcopyright{} \emph{Harvard Math 22a}	
}
\renewcommand{\headrulewidth}{0pt}

%\fancyhead[C]{}
%	\chead{}	
%	\lhead{}
%	\rhead{}
%	\cfoot{}


% Putting copyright on the footer of each page
%\fancypagestyle{secondpage}{
%	\fancyhead[C]{TESTing again} % change this to be blank
%		%\chead{}	
%	\fancyhead[L]{bears} % change this to be blank
%	\fancyhead[R]{dogs} % change this to be blank
%	\fancyfoot[R]{ %keeps the footer the same
%		\vspace*{\fill}\footnotesize\textcopyright{} \emph{Harvard Math 22a}	
%	}
%}
%\renewcommand{\headrulewidth}{0pt}
%\pagestyle{fancy}
\pagestyle{empty}


%\lhead{}
%\rhead{}
%\rfoot{\vspace*{\fill}\footnotesize\textcopyright{} \emph{Harvard Math 22a}}
%\renewcommand{\headrulewidth}{0pt}
%\pagestyle{fancy}




%%
%% Common shortcut commands:
%%
\newcommand{\ds}{\displaystyle}
\newcommand{\sgn}{\operatorname{sgn}}
\renewcommand{\Pr}{\operatorname{P}}
\newcommand{\CPr}[3][{}]{\Pr_{\text{#1}}\left(#2\,|\,#3\right)}

\newcommand{\C}{\mathbb{C}}
\newcommand{\R}{\mathbb{R}}
\newcommand{\Z}{\mathbb{Z}}
\newcommand{\N}{\mathbb{N}}
\newcommand{\Q}{\mathbb{Q}}
\newcommand{\D}{\mathbb{D}}

\newcommand{\rref}{\operatorname{rref}}
\newcommand{\im}{\operatorname{im}}
\newcommand{\rank}{\operatorname{rank}}
\newcommand{\diag}{\operatorname{diag}}
\newcommand{\Span}{\operatorname{span}}
%\renewcommand{\vec}[1]{\mathbf{#1}}
\newcommand{\charpoly}{\mathscr{P}}
\newcommand{\nicefrac}[2]{\sfrac{#1}{#2}} % using xfrac package
\newcommand{\topology}{\mathcal{T}}
\newcommand{\Inn}{\operatorname{Inn}}

\newcommand{\blankbox}[2]{\fbox{\rule{#1}{0in}\rule{0in}{#2}}}

\newenvironment{myindentpar}[1]%
 {\begin{list}{}%
         {\setlength{\leftmargin}{#1}
          \setlength{\rightmargin}{\leftmargin}}%
         \item[]%
 }
 {\end{list}}
\newtheorem{theorem}{Theorem}
\newtheorem*{NStheorem}{Theorem}
\newtheorem{cor}[theorem]{Corollary}
\newtheorem*{claim}{Claim}
\newtheorem{defn}{Definition}

\newenvironment{amatrix}[1]{%
  \left(\begin{array}{@{}*{#1}{c}|c@{}}
}{%
  \end{array}\right)
}

%--------grstep
% For denoting a Gauss' reduction step.
% Use as: \grstep{\rho_1+\rho_3} or \grstep[2\rho_5 \\ 3\rho_6]{\rho_1+\rho_3}
\newcommand{\grstep}[2][\relax]{%
   \ensuremath{\mathrel{
       {\mathop{\longrightarrow}\limits^{#2\mathstrut}_{
                                     \begin{subarray}{l} #1 \end{subarray}}}}}}
\newcommand{\swap}{\leftrightarrow}


%\usetikzlibrary{arrows}
%\usepackage{wrapfig}

%\makeatletter
%\renewcommand*\env@matrix[1][*\c@MaxMatrixCols c]{%
%	\hskip -\arraycolsep
%	\let\@ifnextchar\new@ifnextchar
%	\array{#1}}
%\makeatother
%
%\def\env@matrix{\hskip -\arraycolsep
%	\let\@ifnextchar\new@ifnextchar
%	\array{*\c@MaxMatrixCols c}}

\section{{Change of Coordinates, $\vec\S$4.6}}% 
\chead{\Large\textbf{Change of Coordinates, $\vec\S$4.6}}% 

%\begin{document}
\thispagestyle{fancy}

Recall from last class:

\begin{itemize}
\item If $V$ is spanned by a finite set, then $V$ is called {\bf finite dimensional,} and the {\bf dimension} of $V$, written $\text{dim}(V)$, is the number of vectors in a basis of $V$. Define $\text{dim}(\{0_V\})=0$. If $V$ is not spanned by any finite set, then $V$ is called {\bf infinite dimensional}.
\item {\bf Theorem:} Let $H$ be a subspace of a finite dimensional vector space $V$. Any linearly independent set in $H$ can be expanded to a basis for $H$, and $H$ is finite dimensional and $\text{dim}(H)\leq \text{dim}(V)$.
\item {\bf Basis Theorem:} Let $V$ be an $n$-dimensional vector space with $n\geq 1$. Any linearly independent set of exactly $n$ vectors in $V$ is a basis for $V$. Any set of exactly $n$ vectors that spans $V$ is a basis for $V$.
\item Let $A$ be an $m\times n$ matrix with rows $a_1, \dots, a_m$. Define the {\bf row space} of $A$ by $$\text{Row}(A)=\text{Span}\{a_1, \dots, a_m\}.$$
\item {\bf Theorem:} If $A$ is row equivalent to $B$, then $\text{Row}(A)=\text{Row}(B)$. Furthermore, if $B$ is in reduced row echelon form, the non-zero rows of $B$ form a basis for $\text{Row}(A)$ and $\text{Row}(B)$.
\item The {\bf rank} of the $m\times n$ matrix $A$ is the dimension of $\text{Col}(A)$. We write $\text{rank}(A)=\text{dim}(\text{Col}(A))$. The {\bf nullity} of $A$ is the dimension of $N(A)$. 
\item {\bf Rank-Nullity Theorem:} Let $A$ be an $m\times n$ matrix.
\begin{enumerate}
\item $$\text{dim} \;\text{Row}(A)= \text{dim} \;\text{Col}(A)$$
\item $$\text{rank(A)}+\text{dim} \;N(A) =n$$
\end{enumerate}
\item {\bf Invertible Matrix Theorem Continued} The following are equivalent:
\begin{itemize}
\item[m.] The columns of $A$ form a basis of $\mathbb{R}^n$.
\item[n.] $\text{Col}(A)=\mathbb{R}^n$.
\item[o.] $\text{dim} \; \text{Col}(A)=n$.
\item[p.] $\text{rank}(A)=n$.
\item[q.] $N(A)=\{0\}$.
\item[r.] $\text{dim}\; N(A) =0$.
\end{itemize}
\item {\bf Theorem:} Let $T: V \to W$ be a linear transformation between vector spaces. If $V$ is finite dimensional, then $$\text{dim} \; \text{im}(T) + \text{dim} \; \text{ker}(T) =\text{dim}(V).$$
\end{itemize}

\newpage

{\bf Rank-Nullity Theorem:} Let $A$ be an $m\times n$ matrix.
\begin{enumerate}
\item[a.] $$\text{dim} \;\text{Row}(A)= \text{dim} \;\text{Col}(A)$$
\item[b.] $$\text{rank(A)}+\text{dim} \;N(A) =n$$
\end{enumerate}

\begin{enumerate}

\Problem 

\Part Let $A$ be a 7 by 5 matrix. What is the largest possible rank?

\vfill

\Part Let $A$ be a 3 by 7 matrix. What is the smallest possible nullity?

\vfill

\Part Let $A$ be a 5 by 6 matrix. Suppose the solutions to $Ax=0$ are all multiple of one non-zero vector. Will the matrix equation $Ax=b$ always be consistent?

\vfill


{\bf Invertible Matrix Theorem Continued} The following are equivalent to $A$ being invertible:
\begin{itemize}
\item[m.] The columns of $A$ form a basis of $\mathbb{R}^n$.
\item[n.] $\text{Col}(A)=\mathbb{R}^n$.
\item[o.] $\text{dim} \; \text{Col}(A)=n$.
\item[p.] $\text{rank}(A)=n$.
\item[q.] $N(A)=\{0\}$.
\item[r.] $\text{dim}\; N(A) =0$.
\end{itemize}


{\bf Theorem:} Let $T: V \to W$ be a linear transformation between vector spaces. If $V$ is finite dimensional, then $$\text{dim} \; \text{im}(T) + \text{dim} \; \text{ker}(T) =\text{dim}(V).$$


\newpage

\Problem Let $\mathcal{B}=\{b_1, b_2\}$ and $\mathcal{C}=\{c_1, c_2\}$ be bases for a vector space $V$ where $b_1= c_1+c_2$ and $b_2=-c_1+c_2$. How can we change from the basis $\mathcal{B}$ to $\mathcal{C}$?

\vfill

{\bf Theorem:} Let $\mathcal{B}=\{b_1, \dots, b_n\}$ and $\mathcal{C}=\{c_1, \dots, c_n\}$ be bases of a vector space $V$. Then there exists a unique $n\times n$ matrix $P$ such that $[x]_{\mathcal{C}}=P [x]_{\mathcal{B}}$.

\Problem Prove the theorem.

\vfill

\newpage

{\bf Note:} We call $P$ the {\bf change of coordinates matrix from $\mathcal{B}$ to $\mathcal{C}$}. In general $$P= [[b_1]_\mathcal{C} \; \dots \; [b_n]_\mathcal{C}].$$ Also $P^{-1}$ is the change of coordinates matrix from $\mathcal{C}$ to $\mathcal{B}$.

\Problem Let $b_1=\begin{bmatrix} 1\\ -3\\ \end{bmatrix}$, $b_2=\begin{bmatrix} -2\\ 4\\ \end{bmatrix}$, $c_1=\begin{bmatrix} -7\\ 9\\ \end{bmatrix}$, and $c_2=\begin{bmatrix} -5\\ 7\\ \end{bmatrix}$. Find the change of coordinates matrix from $\mathcal{B}$ to $\mathcal{C}$.

\vfill

\Problem Let $\mathcal{B}=\{1-2t+t^2, 3-5t+4t^2,2t+3t^2\}$ and $\mathcal{E}=\{1, t, t^2\}$ be bases for $\mathbb{P}_2$. Find the change of coordinates matrix from $\mathcal{B}$ to $\mathcal{E}$.

\vfill

\Problem Find $[-1+2t]_\mathcal{B}$.

\vfill

\newpage

Let $V$ be a vector space with basis $\mathcal{B}=\{b_1, \dots, b_n\}$ and $W$ be a vector space with basis $\mathcal{C}=\{c_1, \dots, c_m\}$. Given a linear transformation $T: V\to W$, we naturally seek an associated matrix to represent $T$ using the bases; namely we want an $m\times n$ matrix $M$ such that $[T(x)]_\mathcal{C}=M [x]_\mathcal{B}$ for all $x \in V$. We call $M$ the {\bf matrix of $T$ relative to bases $\mathcal{B}$ and $\mathcal{C}$}. In general $M= [[T(b_1)]_\mathcal{C} \; \dots \; [T(b_n)]_\mathcal{C}]$.

\Problem Let $\mathcal{B}=\{b_1, b_2, b_3\}$ and $\mathcal{D}=\{d_1, d_2\}$ be bases for $V$ and $W$ respectively. Let $T: V \to W$ be the linear transformation such that $T(b_1)=3d_1-5d_2$, $T(b_2)=-d_1+6d_2$, and $T(b_3)=4d_2$. Find the matrix of $T$ relative to $\mathcal{B}$ and $\mathcal{D}$.

\vfill


\end{enumerate}

\begin{comment}
\newpage
\chead{\Large\textbf{Change of Coordinates -- Answers / Solutions}}%
\lhead{}
\rhead{}
\thispagestyle{fancy}

\begin{enumerate}
  \Answer 
  Since we are given $b_1$ as a linear combination of $c_1$ and $c_2$, we can give the coordinates of $b_1$ in terms of basis $\mathcal{C}$; i.e. $$[b_1]_\mathcal{C}=\begin{bmatrix} 1\\ 1\\ \end{bmatrix}.$$ Similarly, we get $$[b_2]_\mathcal{C}=\begin{bmatrix} -1\\ 1\\ \end{bmatrix}.$$. Let $x \in V$, then $x= c_1 b_1+c_2 b_2$ for some scalars $c_1$ and $c_2$. Thus $$[x]_\mathcal{C}=[c_1 b_1 +c_2 b_2]_\mathcal{C}.$$ Since the coordinate mapping is linear, we get $$[x]_\mathcal{C}=c_1 [b_1]_\mathcal{C} +c_2 [b_2]_\mathcal{C}= \begin{bmatrix} 1 & -1\\ 1 & 1\\ \end{bmatrix} \begin{bmatrix} c_1\\ c_2\\ \end{bmatrix} = \begin{bmatrix} 1 & -1\\ 1 & 1\\ \end{bmatrix}[x]_\mathcal{B}.$$ Thus we can change form basis $\mathcal{B}$ to $\mathcal{C}$ by multiplying by the matrix $P=\begin{bmatrix} 1 & -1\\ 1 & 1\\ \end{bmatrix}$.
    \Answer For the proof, we carry out the strategy in problem 1 more generally. 
    
    \begin{proof}
    Let $x\in V$, then there exists $d_1, \dots, d_n \in \mathbb{R}$ such that $$[x]_\mathcal{B}=\begin{bmatrix} d_1\\ \vdots \\ d_n\\ \end{bmatrix}.$$ Then using the linearity of the coordinate mapping and the definition of matrix multiplication, we get
    \begin{align*}
    [x]_\mathcal{C} & = [d_1 b_1 +\dots+ d_n b_n]_\mathcal{C}\\
    & = d_1 [b_1]_\mathcal{C} +\dots + d_n [b_n]_\mathcal{C}\\
    & = [ [b_1]_\mathcal{C} \; \dots \; [b_n]_\mathcal{C} ] \begin{bmatrix} d_1\\ \vdots\\ d_n \end{bmatrix}\\
    & = P [x]_\mathcal{B}.
    \end{align*}
    Thus the matrix $P$ is the change of coordinates matrix from basis $\mathcal{B}$ to basis $\mathcal{C}$.
    \end{proof}
     \Answer Recall the matrix $P$ that changes coordinates from basis $\mathcal{B}$ to basis $\mathcal{C}$ is given by $$P=[[b_1]_\mathcal{C} ]; [b_2]_\mathcal{C} ].$$ Thus we want to find scalars $a_1, a_2$ such that $b_1=a_1 c_1+a_2 c_2$. Thus we want to solve the matrix equation $[ c_1\; c_2] \begin{bmatrix} a_1\\ a_2\\ \end{bmatrix}= b_1$. We could row reduce the augmented matrix $[c_1\; c_2 | b_1]$ or we could use $\begin{bmatrix} a_1\\ a_2\\ \end{bmatrix} = [c_1 \; c_2 ]^{-1} b_1$. Using the inverse matrix, we get $$\begin{bmatrix} a_1\\ a_2\\ \end{bmatrix}=\begin{bmatrix} -7 & -5\\
    9 & 7\\ \end{bmatrix}^{-1}\begin{bmatrix} 1\\ -3\\ \end{bmatrix}=\frac{1}{-4} \begin{bmatrix} 7 & 5\\ -9 & -7\\ \end{bmatrix} \begin{bmatrix} 1\\ -3\\ \end{bmatrix} =\begin{bmatrix} 2\\ -3\\ \end{bmatrix}.$$ Thus $[b_1]_\mathcal{C}=\begin{bmatrix} 2\\ -3\\ \end{bmatrix}.$
     
     Similarly, we have $b_2 = \begin{bmatrix} -2\\ 4\\ \end{bmatrix} = \begin{bmatrix} -7 & -5\\ 9 &7 \\ \end{bmatrix} \begin{bmatrix} a_1 \\ a_2\\ \end{bmatrix}$. Again solving the matrix equation, we get $\begin{bmatrix} a_1\\ a_2\\ \end{bmatrix} = \begin{bmatrix} \frac{-3}{2} \\ \frac{5}{2}\\ \end{bmatrix}$. Thus $[b_2]_\mathcal{C}=\begin{bmatrix} \frac{-3}{2} \\ \frac{5}{2}\\ \end{bmatrix}$.
     
    Thus $$P=\begin{bmatrix} 
    2 & -\frac{3}{2}\\
    -3 & \frac{5}{2}\\
    \end{bmatrix}.$$
    
    Equivalently we could use that $P$ is equal to first changing coordinates from $\mathcal{B}$ to the standard coordinates (given by $P_\mathcal{B}$) and then changing from the standard coordinates to $\mathcal{C}$ (given by $P_\mathcal{C}^{-1}$). Thus $$P=P_\mathcal{C}^{-1} P_\mathcal{B}=\begin{bmatrix} -7 & -5\\
    9 & 7\\ \end{bmatrix}^{-1} \begin{bmatrix} 1 & -2 \\ -3 & 4\\ \end{bmatrix} =\begin{bmatrix} 
    2 & -\frac{3}{2}\\
    -3 & \frac{5}{2}\\
    \end{bmatrix}.$$
    \Answer Let $P$ denote the change of coordinates matrix from $\mathcal{B}$ to the standard coordinates, then $$P=\begin{bmatrix} 1 & 3 & 0\\ -2 & -5 & 2\\ 1 & 4 & 3\\ \end{bmatrix}.$$
     \Answer Recall $$[x]_\mathcal{E}= P [x]_\mathcal{B}.$$ We want to solve the matrix equation $P[x]_\mathcal{B}=\begin{bmatrix} -1\\ 2\\ 0\\ \end{bmatrix}$. Solving the augmented matrix we get
     $$\begin{amatrix}{3}
     1 & 3 & 0 & -1\\
     -2 & -5 & 2 & 2\\
     1 & 4 & 3 & 0\\
     \end{amatrix}\sim \begin{amatrix}{3}
     1 & 0 & 0 & 5\\
     0 & 1 & 0 & -2\\
     0 & 0 & 1 & 1\\
     \end{amatrix}.$$
     Thus $[-1+2t]_\mathcal{B} = \begin{bmatrix} 5\\ -2\\ 1\\ \end{bmatrix}$, and hence $$x=5(1-2t+t^2)-2(3-5t+4t^2)+1\cdot (2t +3t^2).$$
     
       \Answer Recall the matrix $M$ of $T$ relative to bases $\mathcal{B}$ and $\mathcal{D}$ is given by $M=[[T(b_1)]_\mathcal{D} \; [T(b_2)]_\mathcal{D} \; [T(b_3)]_\mathcal{D} ]$. Since $T(b_1)= 3d_1-5d_2$, we get $[T(b_1)]_\mathcal{D}=\begin{bmatrix} 3\\ -5\\ \end{bmatrix}$. Since $T(b_2)= -d_1+6d_2$, we get $[T(b_2)]_\mathcal{D}=\begin{bmatrix} -1\\ 6\\ \end{bmatrix}$. Since $T(b_3)= 4d_2$, we get $[T(b_3)]_\mathcal{D}=\begin{bmatrix} 0\\ 3\\ \end{bmatrix}$. Thus the matrix of $T$ relative to bases $\mathcal{B}$ and $\mathcal{D}$ is given by $\begin{bmatrix} 3 & -1 & 0\\ -5  & 6 & 4\\ \end{bmatrix}.$

\end{enumerate}  
\end{comment}

%\end{document}


%%% Local ables: 
%%% mode: latex
%%% TeX-master: t
%%% End: 
